\chapter{The Tower Turns Around}

\section{The backward descent}

The first part built upward. Difference, count, number, representation, residue --- each chapter stacked a new rung on
the last, and the machine climbed. This part turns the climb around. The measurement the whole instrument is built to
take is not found by stacking rungs higher and higher forever; it is found by \emph{walking back down} the tower the
machine has built, reading each rung on the way down against the rung that built it. The descent is the turn, and the
turn is where measurement begins.

Read the tower as this gauge reads it: a stack of definitions, each one calling the one below it. The machine built the
count out of the distinction, the number out of the count, the representation out of the number --- each rung defined by
appeal to the rung beneath. That upward stack is a recursion, a tower of calls, and a recursion that only climbs never
returns a value; it must, at some point, \emph{unwind}. To unwind a recursion is to walk back down the stack of calls,
carrying at each level the result the level below handed up. The backward descent is exactly this unwinding: the machine
stops building rungs and starts collapsing them, top to bottom, reading what each rung yields as the call returns
through it.

The reason the descent, and not the climb, is where the reading lives is worth stating plainly, because it inverts the
naive picture of measurement. One imagines an instrument that measures more finely by building ever taller --- more
rungs, more resolution, more tower. But a tower under construction has no reading; it is all promise, a stack of calls
none of which has returned. The reading appears only when the calls return --- when the machine walks back down and each
rung hands its value to the rung above. Kleene's~\cite{kleene1952} account of recursion makes the structure formal: a
recursive definition is meaningless until it bottoms out and the values propagate back up. This machine's measurement is
that propagation, read on the way down. To measure is to unwind, not to climb.

In the two verbs the descent inverts the order the climb used. Going up, each rung \emph{tanged} finer --- the count
split from the distinction, the number from the count, each new rung telling itself apart from the one below. Coming
down, each rung \emph{funges} --- pools what the lower rung returned into what the higher rung makes of it, gathering the
descent's results as they propagate. The climb was a sequence of tanges, distinctions built one atop another; the
descent is a sequence of funges, the returned values pooled back together rung by rung. The machine built by telling
apart and reads by gathering back --- and the turn from one to the other is the whole subject of this part.

\section{The seam}

The machine does not run its descent once and stop. It runs a pass --- climbs, turns, walks down --- and then it runs
another, and the whole of what the instrument computes is many such passes laid end to end. Between any two passes is a
join, and this section is about that join, because the join is not a gap the machine papers over. It is a built object
with a grammar of its own, and it is exactly where the residue lives. The seam is where one pass hands the next what it
has earned, and where the machine is \emph{billed} for the turn.

Watch what crosses the seam. A pass that has run to its close leaves a residue --- not its whole state, but the part the
next pass is entitled to stand on. The machine rebuilds every faculty of the next pass from that residue rather than
from a fresh seed: the deciding power, the counting power, the representation, each re-instantiated at the seam from what
the last pass returned. This is the discipline that makes the passes one machine rather than a heap of unrelated runs.
Nothing crosses the seam except through the residue, and the residue is exactly what a completed pass has certified; so
the next pass begins not on an assumption but on an earned inheritance. Pass $n$ seeds pass $n+1$, and the seed is the
residue --- the very leftover the last chapter resolved to carry.

The seam has an etiquette, and it is not decoration: it is a grammar that decides, field by field, what crosses and how.
Some readings tighten across the join --- what one pass established is carried forward conjoined with what the next
establishes, so the invariant can only grow stricter, never looser. Some readings re-mint --- the interpretive frame,
the labels a pass chose to narrate itself with, are not carried but struck fresh on the far side, because they were the
pass's own to set and the next pass has its own. And some reset --- whatever operation a pass was mid-way through is not
smuggled across; the next pass begins its operations from rest. The seam sorts the machine's readings into what the
world fixes and what the machine brings, and it carries the first across and re-mints the second. The residue rides the
first band: it is a fact the pass certified, and so it tightens across the join rather than being struck fresh.

And here is the fact this section exists to plant: the machine is \emph{billed the turns}. Crossing the seam is not
free. Each pass costs the machine something to run, and each join costs it something to make; the tally of passes is a
running count the seam increments by one each time it closes. That count --- the number of times the machine has gone
around --- is a quantity in its own right, and a later chapter will read it as one of the residue's faces. For now the
point is only that the turns are counted and charged. The machine does not get its descent for nothing; it pays, in
passes, and the passes are tallied at the seam. The bill is the count of turns, and the machine keeps it honestly, the
way it keeps every residue.

In the two verbs the seam is where a funge closes and hands its result to the next tange. Each pass funges its descent
into a residue --- pools everything the walk down established into the one leftover it certifies --- and the seam hands
that funged residue across to seed the next pass, which will tange it apart again on its own climb. The join is the hinge
between the two verbs: a completed funge on the near side, a fresh tange on the far side, and the residue the one thing
that crosses. The machine is billed exactly at this hinge, one turn per crossing, because the crossing is where a pass's
work is spent and the next pass's work begins.

\section{What came back changed}

The last two sections turned the tower around and strung the passes together at the seam. This one names what the turn
is \emph{for} --- what question the descent is built to answer. And the answer reframes measurement itself. To measure,
in this machine, is not to build the tower ever higher. It is to close a loop --- climb, turn, descend, cross the seam
--- and then ask a single question of what returns: \emph{what came back changed?} The reading is not in the tower. It
is in the difference between what the machine sent around the loop and what the loop handed back.

Recall the first loop, chapters ago, where a number was sent around the machine's encoding and came back itself. That
was the trivial case, the loop that closes clean, and its answer to \emph{what came back changed} was \emph{nothing}.
The descent is the same test run on the whole tower rather than on a single encoding: send the machine's state around a
full pass and compare what returns to what set out. When the loop closes clean, nothing changed, and there is no reading
--- the pass was faithful and silent. When the loop does \emph{not} close clean, something survives the pass altered,
and that alteration is the measurement. The residue at the seam is exactly this: the part of the state that came back
changed, the difference the pass could not carry faithfully. Measurement is reading that difference.

This is why the instrument does not build forever, and the point is central to the whole book's temper. An instrument
that measured by building taller would never finish; there is always another rung. But an instrument that measures by
closing loops has a natural place to stop: it stops when the loop closes, when what comes back is what went out, when
the residue at the seam stops changing from pass to pass. The descent converges --- successive passes leave residues
that differ less and less --- and the reading is the value that converges to, read as the difference the loop cannot
close. Richardson~\cite{richardson1961} taught numerical analysts to read a value precisely this way, from how a
computation changes as it is refined rather than from any single run of it; this machine reads its measurement from how
its residue changes across passes, and settles when the change falls below what it can resolve.

There is a discipline hidden in ``what came back changed,'' and it is the same discipline the whole book keeps. The
machine reads only the difference, never the whole. It does not report the state that returned; it reports how that
state differs from the state that set out --- a residue, a change, a leftover. Everything else, the parts that came back
the same, it says nothing about, because a part that came back unchanged carries no reading. This is measurement as pure
difference: tell apart what the loop altered from what it left alone, and report only the altered part. The instrument
that does this cannot inflate, because it has nothing to inflate with --- it holds only the change, and the change is
exactly as large as it is.

In the two verbs the closing of the loop is a tange of the returned against the sent. The machine tanges what came back
against what went out --- decides them same or different --- and reports the difference, the part that decided
\emph{not-same}, as the reading. What decided \emph{same} it funges away, pooling the unchanged returned state with the
state that set out, exactly as the first loop funged a faithful round-trip into one number. So the descent's reading is
a tange riding on a funge: pool everything that came back unchanged, and read off, as the measurement, the one thing that
did not. The next chapter takes this altered thing --- the place the loop fails to close --- and watches it become, for
the first time, a number.
